\documentclass[11pt]{article}

\usepackage[margin=1in]{geometry}
\usepackage[T1]{fontenc}
\usepackage[utf8]{inputenc}
\usepackage{booktabs}
\usepackage{longtable}
\usepackage{array}
\usepackage{xcolor}
\usepackage{hyperref}
\usepackage{listings}
\usepackage{amsmath}
\usepackage{siunitx}

\hypersetup{
  hidelinks,
  pdftitle={MORPHEUS Manual},
  pdfauthor={Merlino Development Team}
}

\definecolor{MerlinoBlue}{HTML}{245C73}
\definecolor{MerlinoSlate}{HTML}{4A5561}
\definecolor{MerlinoGold}{HTML}{C88A2D}
\definecolor{MerlinoBg}{HTML}{F7F8F6}
\emergencystretch=2em

\lstdefinestyle{merlino}{
  basicstyle=\ttfamily\small,
  breaklines=true,
  columns=fullflexible,
  keepspaces=true,
  frame=single,
  rulecolor=\color{MerlinoSlate!45},
  backgroundcolor=\color{MerlinoBg},
  xleftmargin=0.5em,
  xrightmargin=0.5em
}

\newcommand{\program}{\texttt{MORPHEUS}}
\newcommand{\merlino}{\texttt{Merlino}}
\newcommand{\xyzin}{\texttt{xyzin}}
\newcommand{\gicforge}{\texttt{GICForge}}
\newcommand{\angstrom}{\text{\AA}}

\title{\textbf{MORPHEUS Manual}\\
\large Algorithmic model generation and semiexperimental refinement in Merlino}
\author{Merlino Development Notes}
\date{\today}

\begin{document}
\maketitle

\begin{abstract}
\program{} is the \merlino{} program for semiexperimental equilibrium-structure
refinement from rotational constants of isotopologues, vibrational and
electronic corrections, and optional reference structural information.  The
program fits Cartesian molecular structures while automatically constructing a
statistically defined active model from topology, symmetry, isotopic coverage,
constraints, predicates and parameter classes.  This manual describes the
recommended command-line and graphical workflows, the input formats, the
constraint language, the numerical options and the standard output files.
\end{abstract}

\tableofcontents
\newpage

\section{Purpose and Design Principles}

\program{} refines an equilibrium geometry by minimizing the difference between
corrected experimental rotational information and the values calculated from a
current Cartesian structure.  The recommended observable is the principal
moment of inertia, because it avoids the reciprocal nonlinearity of rotational
constants and gives a more stable least-squares problem.  Rotational constants
are still reported at the end of every run for direct spectroscopic inspection.

The program is designed around three rules.

\begin{enumerate}
  \item The active model must start from Cartesian atoms and perceived topology,
        not from a user-designed Z-matrix.
  \item The fitted space must be non-redundant, symmetry-adapted and
        statistically well defined.
  \item Constraints and parameter classes should be generated or expressed in
        chemically transparent primitive coordinates, even when the active
        least-squares variables are delocalized GICs or symmetry-adapted
        Cartesian displacements.
\end{enumerate}

Older refinement workflows usually automate the solution of a least-squares
problem after the user has chosen the coordinate model and constraints.
\program{} also automates the formulation of that problem.  It analyzes the
available isotopologues, detects weakly determined directions, preserves the
maximum compatible symmetry and can propose explicit, rank-consistent
constraints or shared parameter classes supported by the data.

\section{Theoretical Background}

\program{} solves a constrained nonlinear least-squares problem in which the
unknowns are not arbitrary Cartesian components but an algorithmically generated
active structural model.  At each geometry \(x\), the calculated rotational
observables are obtained from the principal moments of inertia,
\[
  I_\alpha(x) = \sum_i m_i\left(|\mathbf{r}_i|^2-r_{i\alpha}^2\right),
\]
with isotope masses changed according to each isotopologue.  The recommended
fit observable is the corrected moment of inertia rather than the rotational
constant itself, because it avoids the reciprocal transformation
\(B_\alpha=h/(8\pi^2 c I_\alpha)\) inside the weighted residual.

For an observed rotational constant \(B^0_\alpha\), the target equilibrium
quantity is formed by applying vibrational and electronic corrections with the
convention stored in \xyzin{}:
\[
  B^e_\alpha \simeq B^0_\alpha - \Delta B^\mathrm{vib}_\alpha
                         - \Delta B^\mathrm{el}_\alpha .
\]
The corresponding target moments define the residual vector \(r(x)\).  The
least-squares objective is
\[
  \chi^2(x)=r(x)^T W r(x) + c(x)^T W_c c(x),
\]
where \(W\) contains spectroscopic weights and \(c(x)\) contains optional hard
or soft structural information.  Hard constraints are enforced by projection or
large penalty weights; soft priors enter as additional weighted residuals.

The active displacement is written as
\[
  \Delta x = P\,\Delta y,
\]
where \(P\) is the projector from active model variables to Cartesian
displacements.  For the GIC model, \(P\) is obtained from the frozen
non-redundant, symmetry-labelled GIC B matrix supplied by \gicforge{}.  For the
symmetry-adapted Cartesian model, \(P\) is the translationally and rotationally
cleaned Cartesian symmetry basis.  In both cases, only the totally symmetric
subspace is used for a symmetry-preserving equilibrium refinement.

Primitive constraints, predicates and parameter classes are not separate
coordinate systems.  They are primitive structural relations that are projected
onto the active model.  If \(s(x)\) is a primitive relation, its linearized row
in active variables is
\[
  \frac{\partial s}{\partial y}
  =
  \frac{\partial s}{\partial x}P .
\]
This is the key design rule: users, advisors and imported literature data can
refer to chemically transparent distances, angles, dihedrals and out-of-plane
coordinates, while the solver operates in a rank-consistent active space.

The numerical solver uses weighted normal equations with rank diagnostics,
line-search safeguards and topology checks.  The topology perceived at the
initial geometry is frozen during an ordinary fit.  If a step would create
spurious contacts or change the covalent graph, the step is rejected rather
than allowing the model itself to change during refinement.

\section{Installation and Environment}

The standard project setup is:

\begin{lstlisting}[style=merlino]
merlino-set
merlino-run
\end{lstlisting}

\noindent
The setup script prepares the Python environment used by the GUI and command
line, including RDKit and PySide6 when available.  The Fortran backends are
kept as independent executables and are selected explicitly from the GUI or
with the \texttt{--backend} option.  Check the resolved configuration with:

\begin{lstlisting}[style=merlino]
python -m merlino config
python -m merlino backends
\end{lstlisting}

\noindent
The current production backend for MORPHEUS/SEfit is Python.  The Fortran77 backend
is maintained where required for common Merlino services, especially
\gicforge{} consistency, but the Python MORPHEUS implementation is the
reference path for the full adaptive least-squares solver and reports.

\section{Overall Workflow}

A reproducible \program{} calculation has five stages.

\begin{enumerate}
  \item Prepare a parent Cartesian geometry and isotopologue data.
  \item Convert external sources into the common \xyzin{} molecular container.
  \item Validate the canonical input and inspect the generated diagnostics.
  \item Inspect active coordinates, constraints, ranks, weights and warnings.
  \item Run the fit and inspect residuals, uncertainties, correlations and the
        optimized-space Hessian.
\end{enumerate}

External files are import sources only.  Once a \xyzin{} file exists, all
\merlino{} modules should read geometry, symmetry, rotational, vibrational and
isotopologue information from that single container.  \program{} can update the
container during preprocessing, but the fit itself rereads the canonical data.

\section{Coordinate Models}

\subsection{GIC model}

The default active model is:

\begin{lstlisting}[style=merlino]
coordinate_model = "gic"
\end{lstlisting}

\gicforge{} reads atomic numbers and Cartesian coordinates, builds primitive
internal coordinates, reduces them to a non-redundant generalized internal
coordinate (GIC) set, assigns symmetry labels and provides the active
coordinates and B matrix.  The primitive classes include bond stretches,
valence bends, dihedrals, out-of-plane coordinates, linear bends and ring
coordinates.  Ring coordinates are based on endocyclic dihedrals and analogous
valence-angle combinations rather than Cremer--Pople variables.

Only totally symmetric coordinates are optimized in a symmetry-preserving
refinement.  The Gaussian-style input writer can therefore place the totally
symmetric GICs first and all other symmetry species after the blank separator
when an external geometry optimization should preserve symmetry.

\subsection{Symmetry-adapted Cartesian model}

The alternative active model is:

\begin{lstlisting}[style=merlino]
coordinate_model = "cartesian_symmetry"
\end{lstlisting}

\noindent
This model removes translations and rotations from Cartesian displacements and
symmetry-adapts the remaining space.  Only totally symmetric Cartesian
directions are fitted.  This option does not require an active GIC B matrix,
but constraints are still defined in primitive internal coordinates and are
projected onto the Cartesian active space.

\subsection{Choosing a model}

Use the GIC model for production structural refinements when chemically local
constraints, parameter classes and primitive uncertainties are central to the
interpretation.  Use the symmetry-adapted Cartesian model as an independent
coordinate check or when a fully democratic Cartesian active space is preferred.
Both models return final Cartesian coordinates, from which all primitive
geometrical parameters and their uncertainties are propagated.

\section{Input Files}

\subsection{Canonical \texttt{xyzin} container}

\xyzin{} is the common \merlino{} molecular container.  It starts with a normal
XYZ block and then stores optional uppercase sections.  The section used by
\program{} is \texttt{\#ISOTOPOLOGUES}.

\begin{lstlisting}[style=merlino]
#ISOTOPOLOGUES
SCHEMA merlino.xyzin.isotopologues.v1
UNITS ROTATIONAL=MHz DELTAVIB=MHz DELTAEL=MHz SIGMA=MHz
INDEXING ATOMS=ONE_BASED
BEGIN parent
DEFINITION parent
ROTATIONAL_MHZ A=1000.000 B=800.000 C=600.000
DELTAVIB_MHZ A=1.0 B=2.0 C=3.0 SOURCE=qm CONVENTION=subtract
DELTAEL_MHZ A=0.0 B=0.0 C=0.0 SOURCE=none CONVENTION=subtract
SIGMA_MHZ A=0.01 B=0.01 C=0.01
END
BEGIN C1_13
DEFINITION 1:13
ROTATIONAL_MHZ A=990.000 B=790.000 C=590.000
END
\end{lstlisting}

\noindent
\texttt{DEFINITION} uses one-based atom indexes.  Multiple substitutions are
separated by semicolons, for example \texttt{2:2;5:13}.  \texttt{parent} means
no isotopic substitution.  \texttt{SIGMA\_MHZ=inf} is accepted for a component
with zero least-squares weight; each isotopologue must still retain at least
one positive finite component weight.

Validate a container with:

\begin{lstlisting}[style=merlino]
python -m merlino xyzin-validate --xyzin working/xyzin
python -m merlino xyzin-validate --xyzin working/xyzin --require-semiexp-observations
\end{lstlisting}

\subsection{Standalone job file}

The recommended standalone input extension is \texttt{.mse.toml}.  The schema
identifier is:
\begin{lstlisting}[style=merlino]
merlino.semiexp.job.v1
\end{lstlisting}

\begin{lstlisting}[style=merlino]
schema = "merlino.semiexp.job.v1"
title = "example MORPHEUS fit"

[fit]
backend = "python"
coordinate_model = "gic"
observable = "moments"
rotational_components = "auto"
robust_loss = "none"

[geometry]
units = "angstrom"
atoms = [
  ["O",  0.000000,  0.000000,  0.000000],
  ["H",  0.000000,  0.000000,  0.960000],
  ["H",  0.920000,  0.000000, -0.240000],
]

[[isotopologues]]
label = "parent"
[isotopologues.definition]
substitutions = ""
[isotopologues.constants]
A_MHz = 1000.0
B_MHz = 800.0
C_MHz = 600.0
\end{lstlisting}

Inline \texttt{[[isotopologues]]} records can be replaced by a separate
observations file referenced from a \texttt{[files]} table.  Relative paths are
resolved from the job-file directory.

\subsection{Legacy MSR import}

Legacy MSR files are accepted for compatibility with explicit extensions:
\texttt{.msr}, \texttt{.msr.inp} and \texttt{*\_msr.inp}.  Generic
\texttt{.inp} files are not auto-detected as MSR files.

The importer accepts Z-matrix or Cartesian parent geometries, isotope mass
blocks, \texttt{bexp}, \texttt{dbvib}, \texttt{dbele}/\texttt{dbelec},
\texttt{weights}, frozen Z-matrix variables marked by \texttt{\#}, and
\texttt{PREDICATES} records.  The imported geometry is converted to Cartesians.
Frozen Z-matrix variables are translated to primitive constraints when they
refer to real atoms.  Legacy \texttt{constraints:} records are preserved as
diagnostic compatibility metadata unless they are expressed in the current MORPHEUS
constraint language.

\section{Constraints, Predicates and Parameter Classes}

\subsection{Hard constraints}

Hard constraints remove directions from the least-squares space.  They can be
specified as primitive freezes, Gaussian ModRedundant lines or Gaussian-style
GIC expressions.

\begin{lstlisting}[style=merlino]
R(1,2) Frozen
A 2 1 3 F
DRCH(Frozen,Value=0.0)=R(5,8)-R(5,7)
ROH(Frozen,Value=0.965159)=R(6,7)
\end{lstlisting}

\noindent
The \texttt{Value=} keyword follows Gaussian convention and imposes an explicit
target.  Without \texttt{Value=}, the initial value is frozen.  Constraints are
expanded over symmetry-equivalent primitive parameters when required.

\subsection{Predicates}

Predicates are weighted structural observations rather than exact constraints.
They are useful when accurate electronic-structure or library information
should stabilize an underdetermined refinement without being treated as exact.

\begin{lstlisting}[style=merlino]
python -m merlino semiexp --job input.mse.toml --outdir run \
  --qm-predicate "R(2,3):0.9683:0.002:reference"
\end{lstlisting}

\noindent
The format is \texttt{label\_pattern:value:sigma:source}.  Smaller
\texttt{sigma} values make the predicate stronger.

\subsection{Parameter classes}

Parameter classes tie several active corrections together.  A shared class
does not force final primitive values to be identical; it imposes a common
correction in the active least-squares space.  This is appropriate when
related parameters are not independently determined by the available
isotopologues.

\begin{lstlisting}[style=merlino]
python -m merlino semiexp --job input.mse.toml --outdir run \
  --parameter-class "CH_stretches:shared:R(1,6)|R(2,7)|R(3,8)"
\end{lstlisting}

\noindent
Automatic class suggestions are conservative.  They are generated from
synthon descriptors and then validated in the active totally symmetric GIC
space.  A primitive pattern is accepted only when it is a dominant contribution
to the matched GIC; weak appearances inside highly delocalized GICs are
ignored.  This prevents a class from accidentally capturing unrelated
totally symmetric directions.

\subsection{Model-layer-generated constraints}

The advisor analyzes symmetry, isotopic coverage and reference structural
information.  Typical actions are:

\begin{itemize}
  \item freeze hydrogen local geometry when corresponding deuterations are
        absent;
  \item preserve near-\(C_s\) or higher symmetry by constraining coordinates
        that should vanish in the higher-symmetry limit;
  \item generate shared classes for chemically equivalent corrections based
        on synthon descriptors;
  \item report weak or poorly observed directions before the fit.
\end{itemize}

\begin{lstlisting}[style=merlino]
python -m merlino semiexp --job input.mse.toml --outdir run \
  --fix-hydrogens \
  --qm-predicate "R(2,3):0.9683:0.002:reference" \
  --parameter-class "CH_stretches:shared:R(1,6)|R(2,7)|R(3,8)"
\end{lstlisting}

\section{Least-Squares Solver}

\program{} uses an adaptive Levenberg--Marquardt trust-region solver with
rank-revealing SVD steps, dynamic column scaling, optional robust losses,
line-search safeguards, checkpoint/restart support and deterministic
conditioning diagnostics.  The GIC coordinate definition is built once for a
fresh fit and then reused; the B matrix is recomputed or updated when needed.
Restart jobs rebuild the GIC definition deterministically from the restart
geometry.

\begin{longtable}{@{}p{0.30\textwidth}p{0.60\textwidth}@{}}
\toprule
Option & Meaning \\
\midrule
\endhead
\texttt{backend} & Requested numerical backend: \texttt{python} or
\texttt{fortran77}. \\
\texttt{coordinate\_model} & \texttt{gic} or
\texttt{cartesian\_symmetry}. \\
\texttt{observable} & \texttt{moments}, \texttt{rotational\_constants} or
\texttt{auto}.  \texttt{moments} is recommended. \\
\texttt{rotational\_components} & \texttt{auto}, \texttt{ABC},
\texttt{AB}, \texttt{AC} or \texttt{BC}.  Planar molecules use only two
independent components. \\
\texttt{max\_iter} & Explicit iteration cap.  If omitted, the cap is
proportional to the number of effective parameters. \\
\texttt{damping} & Initial Levenberg--Marquardt damping. \\
\texttt{max\_step} & Maximum active-coordinate trust-region step norm. \\
\texttt{prune\_condition} & Optional deterministic weak-parameter pruning
target.  \texttt{0.0} disables pruning. \\
\texttt{robust\_loss} & \texttt{none}, \texttt{huber},
\texttt{soft\_l1} or \texttt{cauchy}. \\
\texttt{robust\_scale} & Robust residual scale in weighted units; \texttt{0}
selects an automatic median-absolute-deviation scale. \\
\texttt{leave\_one\_out} & Runs exact leave-one-isotopologue-out refits after
the final fit. \\
\texttt{checkpoint}/\texttt{restart} & Save or restart from a checkpoint JSON.\\
\bottomrule
\end{longtable}

\section{Planar Molecules}

A planar molecule has only two independent rotational constants.  When
\texttt{rotational\_components = "auto"}, \program{} chooses among
\texttt{AB}, \texttt{AC} and \texttt{BC}.  In moment mode these are mapped to
\((I_a,I_b)\), \((I_a,I_c)\) and \((I_b,I_c)\).  The automatic choice favors
the most stable pair according to the propagated instability of the omitted
component and then uses Jacobian conditioning as a tie-breaker.  The omitted
rotational component is still calculated and reported as a diagnostic.

\section{Command-Line Workflows}

\subsection{Input validation}

\begin{lstlisting}[style=merlino]
python -m merlino xyzin-validate --xyzin working/xyzin
python -m merlino xyzin-validate --xyzin working/xyzin \
  --require-semiexp-observations
\end{lstlisting}

\noindent
The validation commands check the canonical Merlino container before a
production fit.  A full MORPHEUS run then materializes or updates \xyzin{},
constructs the coordinate model, applies requested constraints/classes and
writes rank, conditioning, planar-component choice and warning diagnostics to
the output manifest.

\subsection{Production fit}

\begin{lstlisting}[style=merlino]
python -m merlino semiexp \
  --job input.mse.toml \
  --xyzin working/xyzin \
  --outdir working/morpheus/run01 \
  --backend python \
  --coordinate-model gic \
  --observable moments \
  --rotational-components auto
\end{lstlisting}

\subsection{Class comparison}

To test whether shared classes are justified, run the same input twice: once
without the class and once with the proposed shared or fixed class.

\begin{lstlisting}[style=merlino]
python -m merlino semiexp --job p-EBN.msr.inp --outdir run_no_class
python -m merlino semiexp --job p-EBN.msr.inp --outdir run_classed \
  --parameter-class "example:shared:R(1,2)|R(3,4)"
\end{lstlisting}

\noindent
The comparison should be discussed in terms of rank, residuals and condition
number.  Class-tied uncertainties are conditional on the shared/fixed class
model and should not be used as independent evidence for every primitive
parameter.

\subsection{Explaining GIC definitions}

For debugging or method reporting, \gicforge{} definitions can be exported with
dominant primitive contributions:

\begin{lstlisting}[style=merlino]
python -m merlino gic-define --geometry parent.xyz \
  --out gic_definition.json \
  --explain-out gic_explain.json \
  --explain-csv gic_explain.csv
\end{lstlisting}

\noindent
The explanation lists the final GIC labels, symmetry irreps, coordinate kinds,
atom indexes, number of contributing primitives and the largest primitive
weight fractions.

\section{Graphical Interface}

The GUI uses the same services as the command line.  A typical GUI run is:

\begin{enumerate}
  \item create or load a Merlino project;
  \item import Cartesian geometry or a legacy/job input;
  \item inspect or edit isotopologue definitions, rotational constants,
        vibrational corrections, electronic corrections and standard
        deviations;
  \item choose Python or Fortran77 backend for each computational step;
  \item choose GIC or symmetry-Cartesian active coordinates;
  \item validate the input and inspect warnings;
  \item open the semiexperimental model tab and accept, edit or reject model-layer-generated
        constraints/classes with their reason, confidence, atom list and
        evidence source visible;
  \item run the fit and inspect output tables.
\end{enumerate}

The GUI must not keep a private isotopologue store.  Imports from tables, MSR
files or other sources are written into \xyzin{} and then reread by the MORPHEUS
service.

\section{Output Files}

Each run writes a text report, CSV diagnostics and a manifest.  The standard
output directory contains:

\begin{longtable}{@{}p{0.48\textwidth}p{0.42\textwidth}@{}}
\toprule
File & Content \\
\midrule
\endhead
\texttt{semiexp\_report.txt} & Human-readable summary of method, constraints,
diagnostics, residuals and final geometry. \\
\texttt{semiexp\_report.json} & Versioned machine-readable report with method,
model, provenance, diagnostics, warnings, residuals and final geometry. \\
\texttt{semiexp\_geometry.xyz} & Optimized Cartesian structure. \\
\texttt{semiexp\_geometry\_parameters.csv} & Primitive bond lengths, angles,
dihedrals and propagated standard errors. \\
\texttt{semiexp\_rotational\_constants.csv} & Experimental corrected,
calculated and residual rotational constants. \\
\texttt{semiexp\_residuals.csv} & Least-squares residual rows. \\
\texttt{semiexp\_parameters.csv} & Active coordinate values, uncertainties
and class labels. \\
\texttt{semiexp\_covariance.csv} and \texttt{semiexp\_correlation.csv} &
Statistical covariance and correlation matrices. \\
\texttt{semiexp\_svd\_diagnostics.csv} & Singular values, numerical rank and
conditioning diagnostics. \\
\texttt{semiexp\_uncertainty\_diagnostics.csv} & Warnings for unusually large
or poorly determined uncertainties. \\
\texttt{semiexp\_warnings.csv} & Machine-readable warnings. \\
\texttt{semiexp\_hessian.csv} and
\texttt{semiexp\_hessian\_eigenvalues.csv} & Optimized-space Hessian and
stationary-point check. \\
\texttt{semiexp\_kraitchman.csv} & Kraitchman diagnostic values when
applicable. \\
\texttt{semiexp\_checkpoint.json} & Restart coordinates and metadata. \\
\texttt{semiexp\_manifest.json} & Input, option and output provenance. \\
\bottomrule
\end{longtable}

The final report gives Cartesian coordinates and conventional primitive
coordinates in \(\angstrom\) and degrees.  Their uncertainties are obtained by
propagating the fitted covariance through the final Cartesian geometry.

\section{Diagnostics and Troubleshooting}

\begin{longtable}{@{}p{0.33\textwidth}p{0.57\textwidth}@{}}
\toprule
Symptom or warning & Interpretation and action \\
\midrule
\endhead
\texttt{rank reduced} & The experimental data or constraints do not determine
all active directions.  Inspect singular vectors, add justified predicates,
classes or constraints, or reduce the active space. \\
\texttt{large condition number} & The fit is numerically sensitive.  Try
moment observables, automatic planar-component selection, shared classes or
better experimental uncertainties. \\
\texttt{large parameter uncertainty} & The parameter is weakly observed.
Check isotopic coverage and whether a chemically justified class or predicate
is needed. \\
\texttt{low robust weight} & An isotopologue behaves as an outlier under the
selected robust loss.  Verify its assignment, constants and corrections. \\
\texttt{planar pair warning} & The selected two-component pair is poorly
conditioned.  Use \texttt{rotational\_components=auto} unless there is a
specific spectroscopic reason not to. \\
\texttt{not a minimum} & The optimized-space Hessian has negative curvature.
Check constraints, symmetry, data consistency and the reference geometry. \\
\bottomrule
\end{longtable}

\section{Regression Tests}

The main scientific contract tests are:

\begin{lstlisting}[style=merlino]
python -m pytest gui/tests/test_scientific_contracts.py -q
\end{lstlisting}

\noindent
They cover \xyzin{} validation, legacy import, GIC behavior,
constraints, predicates, parameter classes, planar-component selection,
Kraitchman diagnostics, uncertainty propagation and representative benchmark
fits.  The parameter-class regression verifies that automatic suggestions are
filtered in the active totally symmetric GIC space and do not remove unrelated
directions.

\section{Recommended Practice}

\begin{enumerate}
  \item Use Cartesian geometries as the primary structure input.
  \item Materialize all external data into \xyzin{} before production fitting.
  \item Use moments of inertia and \texttt{rotational\_components=auto}.
  \item Validate \xyzin{} and input observations before a fit.
  \item Treat generated constraints and classes as documented model choices:
        inspect them, keep their provenance, and report them.
  \item For incomplete isotopic coverage, prefer chemically justified
        predicates or shared classes over arbitrary parameter deletion.
  \item Check final rotational residuals, propagated primitive uncertainties,
        high correlations, leave-one-out diagnostics and Hessian eigenvalues.
  \item Keep the output manifest with any published or archived result.
\end{enumerate}

\section{Compatibility Notes}

Legacy MSR inputs are supported to recover existing data sets and to compare
results.  They are not used as active Z-matrix models.  \program{} converts the
imported information into Cartesians, isotopologue records, primitive
constraints, predicates and optional classes, then rebuilds the active model
from \gicforge{} or symmetry-adapted Cartesians.  The goal is not to obtain a
different structure, but to obtain the same physical structure from a unique,
topology-derived, symmetry-consistent and statistically documented model.

\end{document}
